\section{Threat Model}\label{sec:threat-model}

A \thunderbolt channel involves three roles (a sequence of holders, a semi-trusted party, and the Bitcoin Chain as a passive settlement layer) communicating through a shared append-only ledger (the \emph{Thunderbolt Ledger}), all operating under an additive key-sharing invariant. We specify the participants and communication model (\autoref{sec:sm-participants}) and state the security goals (\autoref{sec:sm-goals}).

\subsection{Parties, Communication, and Trust}\label{sec:sm-participants}

A \thunderbolt channel has three participant roles:

\begin{itemize}
  \item The \textbf{Holder} (initially $\mathsf{H}_0$) is the party that currently owns the UTXO off-chain. Over the lifetime of the channel, the holder role is handed from one user to the next: $\mathsf{H}_0 \to \mathsf{H}_1 \to \cdots$, where each $\mathsf{H}_s$ is an independent user.
  \item The \textbf{Semi-trusted party} $\mathsf{C}$ maintains the complementary share of the lock key. $\mathsf{C}$ never owns the UTXO and never learns the holder secret; its sole role is to preserve the fixed-key invariant via the compensating update. In practice, $\mathsf{C}$ is instantiated as a threshold committee $\{\mathsf{C}_1, \ldots, \mathsf{C}_n\}$ to distribute trust (\autoref{sec:vess-committee}).
  \item The \textbf{Bitcoin Chain} plays the role of a passive verifier: it sees exactly two transactions over the channel's lifetime (Setup and Finalize), regardless of how many holder transitions occur in between.
\end{itemize}

All off-chain communication is mediated by the \textbf{Thunderbolt Ledger}, a shared append-only bulletin board visible to all participants. Each UTXO channel is independently bound to the ledger via a distinct context identifier, so multiple UTXOs can be transferred in parallel without interference. The ledger itself is not a trusted party: it provides availability and ordering but cannot forge, suppress, or alter entries. Concretely, any authenticated broadcast medium (a smart-contract log, a federated bulletin board, or a dedicated server with append-only semantics) can instantiate the Thunderbolt Ledger.

We consider a static-corruption adversary $\mathcal{A}$ that may corrupt any subset of holders (past or future) and at most $t - 1$ of the $n$ committee members before the channel opens. The adversary is probabilistic polynomial-time (PPT) and interacts with the protocol in the random-oracle model, which \autoref{sec:security} instantiates as the $(\mathcal{F}_{\mathrm{RO}},\mathcal{F}_{\mathrm{PKI}})$-hybrid. This trust model is identical to that of standard threshold-signing deployments (FROST~\cite{komlo2021frost}, ROAST~\cite{ruffing2022roast}).

\subsection{Security Goals}\label{sec:sm-goals}

\thunderbolt aims to provide three properties:

\begin{enumerate}
  \item \textbf{Ownership security:} Only the current holder, the party who knows the current holder secret $x_s$, can authorize a transfer or participate in channel finalization.
  \item \textbf{Fixed-key integrity:} The on-chain lock key $P$ remains unchanged throughout the channel's lifetime, regardless of how many off-chain transfers occur.
  \item \textbf{Transfer correctness:} After a valid transfer from $\mathsf{H}_s$ to $\mathsf{H}_{s+1}$, the new holder obtains a valid holder secret $x_{s+1}$ and the committee's share $y_{s+1}$ is updated consistently with the fixed-key invariant.
\end{enumerate}

\thunderbolt provides no liveness guarantee (fewer than $t$ responsive members freezes the channel) and does not hide the transfer graph.
